\documentclass[11pt]{article}

\usepackage[margin=1in]{geometry}
\usepackage{amsmath,amssymb,amsthm}
\usepackage{bm}
\usepackage{mathtools}
\usepackage{hyperref}
\usepackage{algorithm}
\usepackage{algpseudocode}
\usepackage{booktabs}

\newcommand{\R}{\mathbb{R}}
\newcommand{\E}{\mathbb{E}}
\newcommand{\Cov}{\mathrm{Cov}}
\newcommand{\Var}{\mathrm{Var}}
\newcommand{\tr}{\mathrm{tr}}
\newcommand{\diag}{\mathrm{diag}}
\DeclareMathOperator*{\argmin}{arg\,min}
\DeclareMathOperator*{\argmax}{arg\,max}

\title{Contrast Marginalization in RECOVAR's Latent Posterior Solver}
\author{RECOVAR Development Notes}
\date{March 2026}

\begin{document}
\maketitle

\begin{abstract}
This note derives the exact contrast-marginalized posterior over latent
coordinates used in RECOVAR's embedding stage.  We describe three inference
modes---no contrast, profile-MAP, and marginalized contrast---and show how a
single eigendecomposition of $\Lambda^{1/2} H \Lambda^{1/2}$ reduces the
entire contrast quadrature to elementwise arithmetic over
$[\text{batch}, \text{contrast}, \text{latent}]$ tensors.  We give the full
set of posterior moments needed by PPCA/EM with latent contrast, and discuss
quadrature design.
\end{abstract}

\tableofcontents

%% ========================================================================
\section{Statistical Model}
\label{sec:model}
%% ========================================================================

Work in whitened image coordinates for a single image (or a shared-contrast
group).  The generative model is:
\begin{align}
  z &\sim \mathcal{N}(0, \Lambda), &
  \Lambda &= \diag(\lambda_1, \dots, \lambda_k), \label{eq:prior} \\
  y \mid z, c &\sim \mathcal{N}\bigl(c(m + Bz),\; I\bigr), \label{eq:likelihood}
\end{align}
where $y \in \R^p$ is the whitened image, $m \in \R^p$ is the whitened
projected mean volume, $B \in \R^{p \times k}$ is the whitened projected
basis, and $c > 0$ is the per-image contrast scalar.

\paragraph{Sufficient statistics.}
All image-space information enters through six inner products computed once
in RECOVAR's phase-1 forward pass:
\begin{equation}
  \label{eq:stats}
  H = B^\top B, \quad
  g = B^\top y, \quad
  h = B^\top m, \quad
  t = y^\top m, \quad
  \nu = m^\top m, \quad
  \|y\|^2 = y^\top y.
\end{equation}
These have shapes
$(k \times k)$, $(k)$, $(k)$, scalar, scalar, scalar
per image (with a leading batch dimension in practice).

%% ========================================================================
\section{Conditional Posterior for Fixed Contrast}
\label{sec:fixed-contrast}
%% ========================================================================

For fixed $c$, the negative log joint (up to constants independent of $z$
and $c$) is
\begin{equation}
  \Phi(z, c)
  = \tfrac{1}{2}\|y - c(m + Bz)\|^2 + \tfrac{1}{2} z^\top \Lambda^{-1} z
    - \log p(c).
\end{equation}
Expanding and collecting terms quadratic in $z$:
\begin{equation}
  \Phi(z, c) = \tfrac{1}{2} z^\top A(c)\, z - q(c)^\top z
               + \tfrac{1}{2} r(c) - \log p(c),
\end{equation}
where
\begin{align}
  A(c) &= \Lambda^{-1} + c^2 H, \label{eq:Ac} \\
  q(c) &= c\,(g - c\,h), \label{eq:qc} \\
  r(c) &= \|y\|^2 - 2c\,t + c^2\,\nu. \label{eq:rc}
\end{align}

The conditional posterior is Gaussian:
\begin{equation}
  \label{eq:conditional}
  p(z \mid y, c) = \mathcal{N}(\mu_c, \Sigma_c),
  \qquad
  \mu_c = A(c)^{-1} q(c),
  \qquad
  \Sigma_c = A(c)^{-1}.
\end{equation}

%% ========================================================================
\section{Three Inference Modes}
\label{sec:modes}
%% ========================================================================

%% -----------------------------------------------------------------------
\subsection{Mode 1: No Contrast ($c = 1$)}
\label{sec:no-contrast}
%% -----------------------------------------------------------------------

Set $c = 1$ deterministically:
\begin{equation}
  A = \Lambda^{-1} + H, \qquad
  q = g - h, \qquad
  \mu = A^{-1} q, \qquad
  \Sigma = A^{-1}.
\end{equation}
This is the fastest path and the correct default when contrast estimation
is disabled.

%% -----------------------------------------------------------------------
\subsection{Mode 2: Profile-MAP Contrast}
\label{sec:profile}
%% -----------------------------------------------------------------------

RECOVAR's existing estimator jointly optimizes over $(z, c)$:
\begin{equation}
  \hat{c} = \argmin_c \min_z \Phi(z, c).
\end{equation}
After analytically minimizing over $z$, the profile objective is
\begin{equation}
  \label{eq:profile}
  \mathcal{J}_{\text{prof}}(c)
  = r(c) - q(c)^\top A(c)^{-1} q(c) - 2\log p(c).
\end{equation}

\textbf{There is no $\log\det A(c)$ term.}  This is because optimizing
$\min_z \Phi(z, c)$ for fixed $c$ yields $\Phi(\mu_c, c)$, which does not
involve the determinant---the determinant only appears when integrating
$z$ out.

%% -----------------------------------------------------------------------
\subsection{Mode 3: Marginalized Contrast}
\label{sec:marginalize}
%% -----------------------------------------------------------------------

For exact contrast marginalization we integrate out $z$:
\begin{equation}
  p(c \mid y) \propto p(c) \int p(y \mid z, c)\, p(z)\, dz.
\end{equation}
The Gaussian integral gives
\begin{equation}
  p(y \mid c) \propto |A(c)|^{-1/2}
  \exp\!\Bigl(-\tfrac{1}{2}\bigl[r(c) - q(c)^\top A(c)^{-1} q(c)\bigr]\Bigr),
\end{equation}
so the collapsed negative log marginal score is
\begin{equation}
  \label{eq:marginal}
  \boxed{
    \mathcal{J}_{\text{marg}}(c)
    = r(c) - q(c)^\top A(c)^{-1} q(c) + \log\det A(c) - 2\log p(c).
  }
\end{equation}

\paragraph{Why $\log\det$ matters.}
Different contrast values $c$ yield different posterior covariance volumes
$|\Sigma_c| = |A(c)|^{-1}$.  When comparing marginal likelihoods across
$c$, this volume factor is essential.  It is:
\begin{itemize}
  \item irrelevant when $c$ is fixed (Section~\ref{sec:no-contrast}),
  \item absent in profile-MAP (Section~\ref{sec:profile}),
  \item required for marginalization (this section).
\end{itemize}
The old RECOVAR code is \emph{not} wrong---it targets a different
estimator.

%% ========================================================================
\section{Quadrature for Contrast Marginalization}
\label{sec:quadrature}
%% ========================================================================

Let $\{(c_j, w_j)\}_{j=1}^C$ be a positive quadrature rule on
$[c_{\min}, c_{\max}]$.  The unnormalized posterior contrast weight is
\begin{equation}
  \tilde{\omega}_j
  = w_j \, p(c_j) \, |A(c_j)|^{-1/2}
    \exp\!\Bigl(-\tfrac{1}{2}\bigl[r(c_j) - q(c_j)^\top A(c_j)^{-1} q(c_j)
    \bigr]\Bigr),
\end{equation}
normalized as $\omega_j = \tilde{\omega}_j / \sum_\ell \tilde{\omega}_\ell$.

\paragraph{Default rule.}
We use \textbf{16-point Gauss--Legendre on $[0, 3]$}:
\begin{itemize}
  \item The integrand is smooth in $c$.
  \item 16 nodes achieve spectral convergence (errors
        $< 10^{-3}$ vs.\ 64-node reference, verified empirically).
  \item Far fewer nodes than a uniform 50-point search grid.
\end{itemize}

\paragraph{Compatibility.}  If the user passes explicit nodes without
weights, trapezoid weights are computed automatically.

\paragraph{Contrast prior.}
If $\sigma_c^2 = \infty$, use a flat prior on $[c_{\min}, c_{\max}]$.
If finite:
\begin{equation}
  \log p(c) = -\frac{(c - \mu_c)^2}{2\sigma_c^2} + \text{const}.
\end{equation}

%% ========================================================================
\section{Spectral Reduction}
\label{sec:spectral}
%% ========================================================================

The key computational insight: the one-parameter family
$A(c) = \Lambda^{-1} + c^2 H$ can be simultaneously diagonalized for
\emph{all} $c$ by a single eigendecomposition.

\subsection{Setup}

Let $L = \Lambda^{1/2}$ and define the whitened Gram matrix
\begin{equation}
  G = L\, H\, L = Q\, \diag(d_1, \dots, d_k)\, Q^\top,
  \qquad d_\ell \geq 0,
\end{equation}
where $Q$ is orthogonal.  Precompute:
\begin{equation}
  \alpha = Q^\top L\, g, \qquad
  \beta = Q^\top L\, h, \qquad
  T = L\, Q.
\end{equation}

\subsection{Per-Node Quantities}

For each contrast node $c_j$, define the spectral inverse and
transformed right-hand side:
\begin{align}
  s_\ell(c) &= \frac{1}{1 + c^2 d_\ell}, \label{eq:s} \\
  v_\ell(c) &= c\,\alpha_\ell - c^2\,\beta_\ell, \label{eq:v} \\
  \rho_\ell(c) &= s_\ell(c)\, v_\ell(c). \label{eq:rho}
\end{align}

Then:
\begin{align}
  \mu_c &= T\, \rho(c), \label{eq:mu-spectral} \\
  q(c)^\top A(c)^{-1} q(c) &= \sum_{\ell=1}^k s_\ell(c)\, v_\ell(c)^2,
  \label{eq:quad-spectral} \\
  \log\det A(c) &= \text{const} + \sum_{\ell=1}^k \log(1 + c^2 d_\ell).
  \label{eq:logdet-spectral}
\end{align}

\paragraph{Complexity.}
After the one-time $O(Bk^3)$ eigendecomposition, the entire contrast
pass over $C$ nodes is $O(BCk)$ elementwise arithmetic and reductions
over $[\text{batch}, C, k]$ arrays.  No $[B, C, k, k]$ tensors are
materialized.

\subsection{Why This Is an Optimization, Not a Requirement}

Without the spectral trick, exact marginalization is still possible by
computing a Cholesky factorization of $A(c_j)$ at each node and
evaluating the solve, log-determinant, and weight.  That is
mathematically correct.

The eigendecomposition avoids $C$ repeated factorizations and is the
best way to evaluate \emph{all} nodes and \emph{all} moments
efficiently on GPU.  Empirically this gives 8--26$\times$ speedup
over per-node Cholesky (Table~\ref{tab:benchmark}).

\begin{table}[h]
\centering
\caption{GPU benchmark: spectral vs.\ direct Cholesky
         (single H100, median of 10 runs after warmup).}
\label{tab:benchmark}
\begin{tabular}{rrr rr r}
\toprule
$B$ & $k$ & $C$ & Spectral (ms) & Cholesky (ms) & Speedup \\
\midrule
1\,000 & 4  & 16  & 0.20 & 1.56  &  8$\times$ \\
1\,000 & 10 & 50  & 0.28 & 7.32  & 26$\times$ \\
5\,000 & 20 & 16  & 2.43 & 32.05 & 13$\times$ \\
10\,000& 10 & 16  & 1.05 & 10.08 & 10$\times$ \\
50\,000& 10 & 128 & 8.11 & ---   & --- \\
\bottomrule
\end{tabular}
\end{table}

%% ========================================================================
\section{Posterior Moments}
\label{sec:moments}
%% ========================================================================

\subsection{Latent Moments}

The marginalized posterior moments are computed as weighted sums over
contrast nodes:
\begin{align}
  \E[z \mid y] &= \sum_j \omega_j\, \mu_j, \label{eq:Ez} \\
  \E[zz^\top \mid y] &= \sum_j \omega_j\, (\Sigma_j + \mu_j \mu_j^\top),
  \label{eq:Ezz} \\
  \Cov[z \mid y] &= \E[zz^\top \mid y] - \E[z \mid y]\,\E[z \mid y]^\top.
  \label{eq:Covz}
\end{align}

In the spectral basis, $\Sigma_j = T\,\diag(s(c_j))\,T^\top$ and
$\mu_j = T\,\rho(c_j)$, so:
\begin{equation}
  \E[zz^\top \mid y]
  = T \Bigl[\diag\Bigl(\sum_j \omega_j\, s(c_j)\Bigr)
    + \sum_j \omega_j\, \rho(c_j)\,\rho(c_j)^\top \Bigr] T^\top.
\end{equation}
The matrix in brackets is $k \times k$ and built from
$[B, C, k]$ reductions---no $[B, C, k, k]$ materialization.

\subsection{Contrast Moments}
\begin{align}
  \E[c \mid y] &= \sum_j \omega_j\, c_j, \\
  \E[c^2 \mid y] &= \sum_j \omega_j\, c_j^2.
\end{align}

\subsection{Cross Moments for PPCA M-Step}
\label{sec:cross-moments}

When the PPCA model includes contrast multiplicatively
\begin{equation}
  \tilde{y}_n(\xi) = C_n(\xi)\, c_n\, \bigl(\mu(\xi) + W(\xi,:)\, z_n\bigr)
  + \varepsilon_n(\xi),
\end{equation}
the M-step normal equations for row $w_\xi = W(\xi,:)$ are
\begin{equation}
  \Bigl[\sum_n |C_n(\xi)|^2\, \E[c_n^2\, z_n z_n^* \mid y_n]
  + \diag(\text{Reg}(\xi))\Bigr] w_\xi^*
  = \sum_n C_n(\xi)^*\, \tilde{y}_n(\xi)\, \E[c_n\, z_n^* \mid y_n]
  - \mu(\xi) \sum_n |C_n(\xi)|^2\, \E[c_n^2\, z_n^* \mid y_n].
\end{equation}
This requires the following contrast-weighted moments from the E-step:
\begin{align}
  \E[c\, z \mid y] &= \sum_j \omega_j\, c_j\, \mu_j,
  \label{eq:Ecz} \\
  \E[c^2\, z \mid y] &= \sum_j \omega_j\, c_j^2\, \mu_j,
  \label{eq:Ec2z} \\
  \E[c^2\, z z^\top \mid y]
  &= \sum_j \omega_j\, c_j^2\, (\Sigma_j + \mu_j \mu_j^\top).
  \label{eq:Ec2zz}
\end{align}

\paragraph{Summary of returned moments.}
The solver returns a structured result with all of the above:
\begin{center}
\begin{tabular}{ll}
\toprule
Field & Expression \\
\midrule
\texttt{mean\_z}              & $\E[z \mid y]$ \\
\texttt{cov\_z}               & $\Cov[z \mid y]$ \\
\texttt{second\_moment\_z}    & $\E[zz^\top \mid y]$ \\
\texttt{mean\_c}              & $\E[c \mid y]$ \\
\texttt{second\_moment\_c}    & $\E[c^2 \mid y]$ \\
\texttt{mean\_cz}             & $\E[cz \mid y]$ \\
\texttt{mean\_c2z}            & $\E[c^2 z \mid y]$ \\
\texttt{second\_moment\_czz}  & $\E[c^2 zz^\top \mid y]$ \\
\texttt{contrast\_weights\_posterior} & $\omega_j$ \\
\bottomrule
\end{tabular}
\end{center}

%% ========================================================================
\section{Numerical Stability}
\label{sec:stability}
%% ========================================================================

\begin{enumerate}
  \item \textbf{Symmetrize $G$} before \texttt{eigh}:
    $G \leftarrow \tfrac{1}{2}(G + G^\top)$.
  \item \textbf{Clip eigenvalues}: $d_\ell \leftarrow \max(d_\ell, 0)$
    to handle small negative values from float32 rounding.
  \item \textbf{Log-determinant via \texttt{log1p}}:
    $\log\det A(c) = \text{const} + \sum_\ell \log(1 + c^2 d_\ell)$
    avoids catastrophic cancellation when $c^2 d_\ell \ll 1$.
  \item \textbf{Logsumexp for weights}:
    $\omega_j = \mathrm{softmax}(\log\tilde{\omega}_j)$ rather than
    direct exponentiation.
  \item \textbf{Symmetrize output}: $\E[zz^\top]$ and $\Cov[z]$ are
    explicitly symmetrized after the $T M T^\top$ transform to correct
    float32 asymmetry.
\end{enumerate}

%% ========================================================================
\section{Scope Restriction}
\label{sec:scope}
%% ========================================================================

The current implementation supports exact marginalization only when there
is \textbf{one contrast scalar per solve}.  This covers:
\begin{itemize}
  \item SPA per-image solves (one $c$ per image),
  \item Grouped shared-label solves with one shared $c$ across the group.
\end{itemize}

The case of one shared $z$ with multiple independent per-image contrasts
inside a group is \textbf{not supported} and will raise
\texttt{NotImplementedError}.

%% ========================================================================
\section{Implementation Map}
\label{sec:code}
%% ========================================================================

\begin{center}
\begin{tabular}{ll}
\toprule
Module / Function & Role \\
\midrule
\texttt{contrast\_posterior.py} & Standalone solver module \\
\quad \texttt{make\_contrast\_quadrature} & Gauss--Legendre / trapezoid / custom \\
\quad \texttt{solve\_no\_contrast} & $c=1$ fast path \\
\quad \texttt{solve\_profile\_contrast} & Profile-MAP (no log-det) \\
\quad \texttt{solve\_marginalized\_contrast} & Full marginalization \\
\quad \texttt{solve\_latent\_posterior} & Dispatch wrapper \\
\midrule
\texttt{embedding.py} & Integration layer \\
\quad \texttt{\_solve\_batch\_from\_stats\_v2} & Routes to solver by mode \\
\quad \texttt{get\_per\_image\_embedding\_multi\_zdim} & Accepts \texttt{contrast\_mode} \\
\bottomrule
\end{tabular}
\end{center}

%% ========================================================================
\appendix
\section{Pseudocode}
\label{app:pseudocode}
%% ========================================================================

\begin{algorithm}[H]
\caption{Marginalized contrast posterior (spectral form)}
\label{alg:marginalized}
\begin{algorithmic}[1]
\Require $H_{b,k,k'},\; g_{b,k},\; h_{b,k},\; t_b,\; \nu_b,\;
  \|y\|^2_b,\; \lambda_k,\; c_j,\; w_j$
\State $L_k \gets \sqrt{\lambda_k}$
\Comment{$[K]$}
\State $G_{b} \gets L \, H_b \, L$; symmetrize
\Comment{$[B,K,K]$}
\State $d_b, Q_b \gets \mathrm{eigh}(G_b)$; $d \gets \max(d, 0)$
\Comment{one batched eigh}
\State $\alpha_b \gets Q_b^\top (L \odot g_b)$,\quad
       $\beta_b \gets Q_b^\top (L \odot h_b)$,\quad
       $T_b \gets L \, Q_b$
\For{each node $c_j$}
  \Comment{vectorized over $[B, C, K]$}
  \State $s_{b,j,\ell} \gets 1/(1 + c_j^2\, d_{b,\ell})$
  \State $v_{b,j,\ell} \gets c_j \alpha_{b,\ell} - c_j^2 \beta_{b,\ell}$
  \State $\rho_{b,j,\ell} \gets s_{b,j,\ell}\, v_{b,j,\ell}$
\EndFor
\State \textit{quad}$_{b,j} \gets \sum_\ell v \cdot \rho$
\State \textit{logdet}$_{b,j} \gets \sum_\ell \log(1 + c_j^2 d_{b,\ell})$
\State $r_{b,j} \gets \|y\|^2_b - 2c_j t_b + c_j^2 \nu_b$
\State $\omega_{b,j} \gets \mathrm{softmax}_j\bigl(
  \log w_j + \log p(c_j) - \tfrac{1}{2}(r - \textit{quad} + \textit{logdet})
  \bigr)$
\State $\E[z]_b \gets T_b \sum_j \omega_{b,j}\, \rho_{b,j}$
\State $M_b \gets \diag(\sum_j \omega_{b,j}\, s_{b,j})
  + \sum_j \omega_{b,j}\, \rho_{b,j}\,\rho_{b,j}^\top$
\Comment{$[B,K,K]$}
\State $\E[zz^\top]_b \gets T_b\, M_b\, T_b^\top$; symmetrize
\State \textit{(analogous for cross moments $\E[cz]$, $\E[c^2 zz^\top]$,
  etc.)}
\end{algorithmic}
\end{algorithm}

\end{document}
